\documentclass[10pt,a4paper]{amsart}
\usepackage[margin=19mm]{geometry}
\usepackage{amsmath,amssymb,mathtools}
\usepackage[scr=rsfs,cal=euler]{mathalfa}
\usepackage{xfrac}
\usepackage[unicode,hidelinks,bookmarks=false]{hyperref}
\newcommand{\ie}{i.e.\ }
\newcommand{\cf}{cf.\ }
\newcommand{\resp}{respectively\ }
\newcommand{\Subsection}{\subsection}
\providecommand{\nobreakdash}{\nobreak\hbox{-}}
\setlength{\emergencystretch}{2em}
\multlinegap0pt
\usepackage{bm}
\usepackage{bbm}
\usepackage{dsfont}
\usepackage{xspace}
\usepackage{cancel}
\usepackage{mathtools}
\usepackage{amsthm}
\makeatletter
\@ifundefined{theorem}{\newtheorem{theorem}{Theorem}}{}
\@ifundefined{lemma}{\newtheorem{lemma}[theorem]{Lemma}}{}
\makeatother
\newcommand{\EE}{\mathbb{E}}
\newcommand{\Var}{\operatorname{Var}}
\renewcommand{\ge}{\geqslant}
\renewcommand{\le}{\leqslant}
\title[An entropy bridge from weighted to spectral Tur\'{a}n theore]{An entropy bridge from weighted to spectral Tur\'{a}n theorems}
\author{Yongtao Li}
\date{Source: arXiv:2608.24388v2 (CC BY 4.0)}
\begin{document}
\maketitle

\noindent\emph{Source and licence.} An entropy bridge from weighted to spectral Tur\'{a}n theorems. Version arXiv:2608.24388v2 (CC BY 4.0), distributed under the Creative Commons Attribution 4.0 International licence, as recorded in the arXiv licence metadata for this exact version and shown on the abstract page https://arxiv.org/abs/2608.24388v2. Full rights notice: licenses/arxiv-2608.24388-CC-BY-4.0.txt.\par\smallskip

\noindent\textit{Editorial scope.} Counted unit: the Theorem whose source label is \texttt{thm:weighted-stability} in the article (source0.tex lines 2250--2266, source label \texttt{thm:weighted-stability}), delivered together with the article's own complete original proof of it (source0.tex lines 2269--2473, one \texttt{proof} environment of 205 lines). No mathematics is written, repaired or re-derived: statement and proof are the article's own text line for line. Required context retained uncounted: thm:removal (lines 925--929); thm:stability (lines 935--947); lem:markov (lines 1187--1194); lem:chebyshev (lines 1198--1206). Every definition, display and label of those lines is retained unchanged. Omitted from this package and not redistributed: 1--924, 930--934, 948--1186, 1195--1197, 1207--2249, 2267--2268, 2474--4490. Nothing inside a retained excerpt is omitted or altered: every retained body is delivered exactly as the article's lines are written, so no omission receipt of any kind accompanies this package. Citations: the retained excerpts carry no citation command at all, so no bibliography entry of the article is transcribed and no reference is redistributed. Reference adaptation: none was needed and none was made; every reference inside the delivered unit resolves inside it, so no unresolved reference or citation is produced. Printed slips kept literal: no printed word, symbol, hypothesis or number of the counted statement or of its proof was altered. Layout and class: the article's own class is \texttt{article} with packages including fontenc, inputenc, amsmath, amssymb, amsthm, mathtools, fullpage, color, hyperref, tikz, bm, geometry, booktabs; the wrapper is the lane's pinned \texttt{amsart} editorial wrapper at 10pt on A4, which restates the theorem-like environments the retained text needs and adds the article's own macro definitions that the retained text needs (\texttt{\textbackslash EE}, \texttt{\textbackslash Var}, \texttt{\textbackslash ge}, \texttt{\textbackslash le}), with extra wrapper packages (bm, bbm, dsfont, xspace, cancel, mathtools). The delivered numbering is the unit's own. Assets: no retained span contains a figure environment, a graphics-inclusion command, a PDF-page inclusion, a table or a TikZ picture; no image file is copied into this package, the package declares no asset dependency and no image file is redistributed with it. Licence: version arXiv:2608.24388v2 of this article is distributed under the Creative Commons Attribution 4.0 International licence, as recorded in the arXiv licence metadata for this exact version and shown on the abstract page https://arxiv.org/abs/2608.24388v2. A full rights notice is delivered at licenses/arxiv-2608.24388-CC-BY-4.0.txt. Characters: every retained excerpt is pure ASCII, so the delivered engine, XeLaTeX with its default Latin Modern text font, reports no missing character.
\begin{theorem}[Graph removal lemma]\label{thm:removal}
Let $F$ be a fixed graph with $f=v(F)$. For every
$\varepsilon >0$, there is $\nu =\nu (F,\varepsilon )>0$
such that every graph $G$ on $N$ vertices containing at most $\nu N^{f}$ copies of $F$ can be made $F$-free by deleting at most $\varepsilon N^{2}$ edges.
\end{theorem} 

\begin{theorem}[Erd\H{o}s--Simonovits stability]
\label{thm:stability}
Let $F$ be a fixed graph with $\chi (F)=r+1\ge 2$. For
every $\varepsilon >0$, there are $\theta =\theta
(F,\varepsilon )>0$ and $N_1=N_1(F,\varepsilon )$ such
that every $F$-free graph $G$ on $N\ge N_1$ vertices with 
$$e(G)\ge \Big(1-\frac1r-\theta \Big)\frac{N^{2}}{2}$$ 
admits a partition $V(G)=B_1\cup \cdots \cup B_r$ with
$\bigl\lvert |B_i|-\frac Nr\bigr\rvert \le \varepsilon N$
for every $i\in [r]$, and with at most $\varepsilon N^{2}$
pairs that lie in different parts and are non-adjacent
in $G$.
\end{theorem}

\begin{lemma}[Markov's inequality]
\label{lem:markov}
Let $Y$ be a nonnegative random variable with finite
mean. Then for every $a>0$,
\[
  \Pr (Y\ge a)\le \frac{\EE [Y] }{a}.
\]
\end{lemma}

\begin{lemma}[Chebyshev's inequality]
\label{lem:chebyshev}
Let $Z$ be a random variable with finite variance.
Then for every $a>0$,
\[
  \Pr \bigl( \big\lvert Z-\EE [Z] \big\rvert \ge a\bigr)
  \le \frac{\Var (Z)}{a^2}.
\]
\end{lemma}

\begin{theorem}[Weighted stability]
\label{thm:weighted-stability}
Let $F$ be a graph with $\chi (F)=r+1\ge 3$. For every
$\varepsilon >0$, there are $\eta =\eta (F,\varepsilon )>0$ and $\delta =\delta(F,\varepsilon)>0$ such that
if $\bm{p}$ is a probability vector with
$\Vert \bm{p} \Vert_\infty\le \eta$, 
and $G$ is an $F$-free graph with $\Phi_G(\bm{p})\ge 1- \frac{1}{r} - \delta$, then there is a
  partition 
  $$V(G)=U_1\cup \cdots \cup U_r$$
   satisfying 
  \[
    \Bigl|p(U_i)-\frac1r\Bigr|\le \varepsilon 
    \quad (\text{for every}~i\in [r]),
    \qquad
    M(G,\bm{p},U_1,\ldots ,U_r)\le \varepsilon .
  \] 
\end{theorem}

\begin{proof}
Let $f:=v(F)$, and assume without loss of generality
that $\varepsilon\le\frac12$. We first fix the
constants. 
Apply Theorem \ref{thm:stability} with parameter
$\varepsilon/8$ to obtain $\theta=\theta(F,\varepsilon/8)>0$
and $N_1=N_1(F,\varepsilon/8)$ such that every $F$-free
graph $D$ on $N\ge N_1$ vertices with
$e(D)\ge\bigl(1-\frac1r-\theta\bigr)\frac{N^2}{2}$
admits a partition $V(D)=B_1\cup\cdots\cup B_r$ with
$\bigl\lvert|B_i|-\frac Nr\bigr\rvert
\le\frac{\varepsilon}{8}N$ for every $i\in[r]$ and
with at most $\frac{\varepsilon}{8}N^2$ pairs lying
in different parts that are non-adjacent in $D$.
Apply Theorem \ref{thm:removal} with parameter
$\theta/8$ to obtain $\nu=\nu(F,\theta/8)>0$ such that
every graph on $N$ vertices containing at most
$\nu N^f$ copies of $F$ can be made $F$-free
by deleting at most $\frac{\theta}{8}N^2$ edges.
Finally, we set  
\[
  \delta:=\frac{\theta}{4},
  \qquad
  \eta:=\min\Bigl\{\frac{\nu}{f^2},\
  \frac{\varepsilon^2}{32r},\
  \frac{\varepsilon}{8}\Bigr\}.
\]



\noindent
\textit{Step 1. Rationalization and blow-up.}
Put $\xi:=\min\{\eta,\frac{\theta}{16},\frac{\varepsilon}{8}\}$. Since $V(G)$ is finite and the rational points are dense, 
we may choose a {\it rational probability vector} $\bm q=(q_v)_{v\in V(G)}$ with
$\sum_{v}\lvert p_v-q_v\rvert\le\xi$ and with $q_v>0$ for every $v\in V(G)$. 
For any set $\mathcal P$ of unordered pairs of distinct vertices, we get 
\[
  2\Bigl\lvert\sum_{uv\in\mathcal P}
  (p_up_v-q_uq_v)\Bigr\rvert
  \le\sum_{u}\sum_{v}
  \bigl(p_u\lvert p_v-q_v \rvert
  +q_v\lvert p_u-q_u \rvert\bigr)\le2\xi .
\]
Taking $\mathcal P=E(G)$, 
and taking $\mathcal P$ to be the set
of cross non-edges of any partition, we obtain 
\begin{equation}\label{eq:rat}
  \lvert\Phi_G(\bm p)-\Phi_G(\bm q)\rvert\le2\xi,
  \quad
  M(G,\bm p, \cdot)\le M(G,\bm q, \cdot)+2\xi,
  \quad
  \lvert p(S)-q(S)\rvert\le\xi
\end{equation}
for every $S\subseteq V(G)$. Moreover
$\lVert\bm q\rVert_\infty\le\eta+\xi\le2\eta
\le\nu/\binom f2$. In particular
\begin{equation}\label{eq:q-density}
  \Phi_G(\bm q)\ \ge\ \Phi_G(\bm p)-2\xi
  \ \ge\ 1-\frac1r-\frac{\theta}{4}
   -\frac{\theta}{8}
  \ \ge\ 1-\frac1r-\frac{3\theta}{8}.
\end{equation}
Let $N_{\bm q}$ be a common denominator of the $q_v$ with $q_v >0$. For a positive integer $t$ put
$N:=tN_{\bm q}$, so that $Nq_v$ is a nonnegative integer
for every $v\in V(G)$. Enlarging $t$ does not change
$\bm q$, so $N$ may be taken as large as we wish, and
we fix $t$ so large that $N\ge N_1$. Construct the
blow-up $B$ of $G$ by replacing each $v\in V(G)$ with an
independent cluster $C_v$ of size $k_v:=Nq_v$ and placing a complete bipartite graph between $C_u$ and
$C_v$ for every $uv\in E(G)$. 
Then $v(B)=N$.  By  \eqref{eq:q-density},
\begin{equation}\label{eq:blowup-density}
  \frac{2e(B)}{N^2}
  =2\sum_{uv\in E(G)}q_uq_v
  =\Phi_{G}(\bm q)
  \ \ge\ 1-\frac1r-\frac{3\theta}{8}.
\end{equation}

\noindent
\textit{Step 2. Using the graph removal lemma and stability.} 
We claim that $B$ contains at most $\nu N^f$  
copies of $F$. Project each vertex of such a copy to
the vertex of $G$ whose cluster contains it. If the
projections were distinct they would form a copy of
$F$ in $G$, contradicting $F$-freeness. Hence two
vertices of the copy lie in a common cluster. For a
fixed pair of vertices of $F$, the number of injective
maps $V(F)\to V(B)$ sending that pair into a common
cluster is at most
\[
  \sum_{v\in V(G)}k_v^2N^{f-2}
  =N^f\sum_{v\in V(G)}q_v^2
  \le N^f\lVert\bm q\rVert_\infty ,
\]
so a union bound over the $\binom f2$ pairs gives
\[
  N_F(B)\le\binom f2N^f\lVert\bm q\rVert_\infty
  \le\nu N^f,
\]
where $N_F(B)$ is the number of copies of $F$
in $B$. By the choice of $\nu$, deleting at most
$\frac{\theta}{8}N^2$ edges from $B$ yields an
$F$-free graph $H$ on $V(B)$.  Then, by
\eqref{eq:blowup-density},
\[
  \frac{2e(H)}{N^2}\ \ge\ \frac{2e(B)}{N^2}
  -\frac{\theta}{4}
  \ \ge\ 1-\frac1r-\theta .
\]
Since $N\ge N_1$, Theorem \ref{thm:stability} applies
to $H$ and yields a partition
$V(B)=B_1\cup\cdots\cup B_r$ such that 
\begin{equation}\label{eq:part-sizes}
  \Bigl\lvert|B_i|-\frac Nr\Bigr\rvert
  \le\frac{\varepsilon}{8}N
\end{equation}
for every $i\in [r]$, and such that at most $\frac{\varepsilon}{8}N^2$ pairs lie
in different parts and are non-adjacent in $H$. As
$E(H)\subseteq E(B)$, every pair non-adjacent in $B$
is non-adjacent in $H$, so
the number of pairs in different parts that are
  non-edges of $B$ is at most $\frac{\varepsilon}{8}N^2$.

\smallskip 
\noindent
\textit{Step 3. Transferring the partition of the blow-up back to $G$.}
The remaining step is to transform the above partition of the blow-up $B$ into a partition of the original graph $G$.  
For every vertex $v\in V(G)$ and every $i\in[r]$, 
we put  
$$q_{vi}:= \frac{| C_v\cap B_i|}{| C_v|}. $$
Then $\sum_{i}q_{vi}=1$. 
Assign each vertex 
$v\in V(G)$
independently to a random set $U_i$ with probability $q_{vi}$. 
We write $\xi_{vi}:=\mathbf 1_{\{v\in U_i\}}$ 
and $Z_i:=q(U_i)=\sum_{v\in V(G)}q_v\,\xi_{vi}$. Since
$\lvert C_v\rvert=Nq_v$, \eqref{eq:part-sizes} gives
\begin{equation}\label{eq:Zi-expectation}
  \EE[Z_i]=\sum_{v\in V(G)}q_v \, q_{vi}
  =\frac1N\sum_{v\in V(G)}\lvert C_v\cap B_i\rvert
  =\frac{|B_i|}{N}
  =\frac1r\pm\frac{\varepsilon}{8}.
\end{equation}
The variables $\xi_{vi}$, with $v\in V(G)$, are independent
with $\Var(\xi_{vi})=q_{vi}(1-q_{vi})$, so
\begin{equation}\label{eq:Zi-variance}
  \Var(Z_i)=\sum_{v\in V(G)}q_v^2 \, 
  q_{vi}(1-q_{vi})
  \le\sum_{v\in V(G)}q_v^2
  \le\lVert\bm q\rVert_\infty\le2\eta ,
\end{equation}
where the last inequality holds because
$\sum_vq_v^2\le \lVert\bm q\rVert_\infty 
\sum_v q_v = \lVert\bm q\rVert_\infty $.

Let $M:=M(G,\bm q, U_1,\ldots,U_r)$. For a
non-edge $uv$ of $G$ the assignment separates $u$ and
$v$ with probability $1-\sum_{i}q_{ui}q_{vi}$, and the
number of pairs of $C_u\times C_v$ lying in different
parts is exactly
$N^2q_uq_v\bigl(1-\sum_{i=1}^r q_{ui}q_{vi}\bigr)$. 
 All these pairs are non-edges of $B$. Hence,
\begin{equation}\label{eq:M-expectation}
  \EE[M]=\sum_{\substack{uv\notin E(G)\\u<v}}
  q_uq_v\Bigl(1-\sum_{i=1}^rq_{ui}q_{vi}\Bigr)
  \le\frac{\#\{\text{cross non-edges of }B\}}{N^2}
  \le\frac{\varepsilon}{8},
\end{equation}
where the inequality holds because the right-hand side also counts pairs of vertices lying in the same cluster $C_v$, which are non-edges of $B$ and have no counterpart on the left. 

By \eqref{eq:Zi-variance} and Chebyshev's inequality (Lemma \ref{lem:chebyshev}), we get that for each $i\in[r]$, 
\[
  \Pr \bigg( \bigl\lvert Z_i-\EE[Z_i]\bigr\rvert
  >\frac{\varepsilon}{2}  \bigg)
  \le\frac{4\Var(Z_i)}{\varepsilon^2}
  \le\frac{8\eta}{\varepsilon^2} \le  \frac{1}{4r}.
\]
Moreover, by \eqref{eq:M-expectation} and Markov's inequality (Lemma \ref{lem:markov}), we get 
$$\Pr\Bigl(M>\frac{\varepsilon}{2}\Bigr)
\le\frac{\varepsilon/8}{\varepsilon/2}=\frac14.$$
Since $r\cdot\frac{1}{4r}+\frac14=\frac12<1$, a union
bound shows that with positive probability none of
these $r+1$ events occurs. Consequently, 
there exists an $r$-partition $V(G)=U_1\cup \cdots \cup U_r$ satisfying 
\[   \Big| Z_i - \mathbb{E}[Z_i]\Big| \le \frac{\varepsilon}{2} \quad 
\text{for every $i\in [r]$}, \quad \text{and} \quad M \le \frac{\varepsilon}{2}.  \]
Recall that $Z_i=q(U_i)$.  By \eqref{eq:Zi-expectation}, the triangle inequality yields 
\[
  \Bigl\lvert q(U_i)-\frac1r\Bigr\rvert
  \le\frac{\varepsilon}{2}+\frac{\varepsilon}{8}
  ~~ \text{for every}~i\in[r],
  \quad \text{and} \quad 
  M(G,\bm q, U_1,\ldots,U_r)\le\frac{\varepsilon}{2}.
\]
Finally, by \eqref{eq:rat} and $\xi\le\varepsilon/8$, using the triangle inequality again, we have 
\[
  \Bigl\lvert p(U_i)-\frac1r\Bigr\rvert
  \le\frac{5\varepsilon}{8}+\xi\le\varepsilon ~~
  \text{for every}~i\in[r], 
  \quad \text{and} \quad 
  M(G,\bm p, U_1,\ldots,U_r)
  \le\frac{\varepsilon}{2}+2\xi\le\varepsilon,
\]
which is the desired assertion.
\end{proof}

\end{document}
